\documentclass[11pt]{article}

\usepackage[T1]{fontenc}
\usepackage[utf8]{inputenc}
\usepackage[spanish,es-nodecimaldot]{babel}
\usepackage{amsmath,amssymb,mathtools}
\usepackage{enumitem}
\usepackage[a4paper,margin=2.5cm]{geometry}
\usepackage[hidelinks]{hyperref}

\newcommand{\Cstar}{C^*}
\newcommand{\GL}{\mathrm{GL}}
\newcommand{\Cb}{C_b}
\newcommand{\Czero}{C_0}
\newcommand{\id}{\mathrm{id}}

\newenvironment{definicion}[1]
  {\par\medskip\noindent\textbf{Definición #1.}\ }
  {\par\medskip}
\newenvironment{teorema}[1]
  {\par\medskip\noindent\textbf{Teorema #1.}\ \itshape}
  {\par\medskip}
\newenvironment{proposicion}[1]
  {\par\medskip\noindent\textbf{Proposición #1.}\ \itshape}
  {\par\medskip}
\newenvironment{corolario}[1]
  {\par\medskip\noindent\textbf{Corolario #1.}\ \itshape}
  {\par\medskip}
\newenvironment{lema}[1]
  {\par\medskip\noindent\textbf{Lema #1.}\ \itshape}
  {\par\medskip}
\newenvironment{observacion}[1]
  {\par\medskip\noindent\textbf{Observación #1.}\ }
  {\par\medskip}
\newenvironment{ejemplo}[1]
  {\par\medskip\noindent\textbf{Ejemplo #1.}\ }
  {\par\medskip}

\title{Enunciados completos de la Parte 1 de $C^*$-álgebras: Teoría espectral en álgebras de Banach}
\author{}
\date{}

\begin{document}
\maketitle

\section*{Índice general}
\begin{itemize}[leftmargin=*]
  \item \textbf{Definiciones:} 1.1, 1.4, 1.7, 1.11, 1.13, 1.17, 1.19, 1.25, 1.27, 1.29.
  \item \textbf{Observaciones:} 1.2, 1.3, 1.9, 1.28.
  \item \textbf{Ejemplos:} 1.5, 1.6, 1.30.
  \item \textbf{Lemas:} 1.14.
  \item \textbf{Proposiciones:} 1.8, 1.12, 1.18, 1.20, 1.21, 1.23, 1.31.
  \item \textbf{Corolarios:} 1.10, 1.24.
  \item \textbf{Teoremas:} 1.15, 1.16, 1.22, 1.26.
\end{itemize}

\section*{Definiciones}


% Fuente: Parte1.pdf, p. 1
\begin{definicion}{1.1}
Un \emph{álgebra de Banach} es un álgebra asociativa $A$ sobre $\mathbb{C}$ con una norma tal que
\begin{itemize}[leftmargin=2em]
  \item $(A,+,\lVert\cdot\rVert)$ es un espacio de Banach;
  \item para todo $x,y\in A$,
  \[
    \lVert xy\rVert\leq \lVert x\rVert\,\lVert y\rVert.
  \]
\end{itemize}
Además decimos que $A$ es:
\begin{itemize}[leftmargin=2em]
  \item \emph{con unidad} si tiene un neutro multiplicativo $1_A$ que satisface
  \[
    \lVert 1_A\rVert=1;
  \]
  \item \emph{conmutativa} si
  \[
    xy=yx
    \qquad\text{para todo }x,y\in A.
  \]
\end{itemize}
\end{definicion}

% Fuente: Parte1.pdf, p. 1
\begin{definicion}{1.4}
Sea $A$ un álgebra compleja. Una \emph{involución} en $A$ es un mapeo
\[
  *:A\longrightarrow A
\]
tal que
\[
  (\lambda a+b)^*=\overline{\lambda}\,a^*+b^*,
  \qquad
  (a^*)^*=a,
  \qquad
  (ab)^*=b^*a^*.
\]
Si $A$ está equipada con una involución, diremos que $A$ es un \emph{álgebra involutiva}.

Si $A$ es además un álgebra de Banach tal que
\[
  \lVert a^*\rVert=\lVert a\rVert,
\]
diremos que $A$ es un \emph{álgebra de Banach involutiva}.
\end{definicion}

% Fuente: Parte1.pdf, p. 2
\begin{definicion}{1.7}
Sea $A$ un álgebra de Banach con unidad. Definimos
\[
  \GL(A)=\{\text{elementos invertibles de }A\}.
\]
\end{definicion}

% Fuente: Parte1.pdf, p. 3
\begin{definicion}{1.11}
Sea $A$ un álgebra de Banach con unidad y $x\in A$. El \emph{espectro} de $x$ se define como
\[
  \sigma(x)=\sigma_A(x)
  =\{\lambda\in\mathbb{C}:\lambda 1_A-x\notin\GL(A)\}.
\]
\end{definicion}

% Fuente: Parte1.pdf, p. 3
\begin{definicion}{1.13}
Sea $A$ un álgebra de Banach y $x\in A$. Definimos el \emph{conjunto resolvente} de $x$ como
\[
  \rho(x)=\mathbb{C}\setminus\sigma(x),
\]
y la \emph{función resolvente}
\[
  R_x:\rho(x)\longrightarrow A,
  \qquad
  R_x(z)=(z1_A-x)^{-1}.
\]
\end{definicion}

% Fuente: Parte1.pdf, p. 4
\begin{definicion}{1.17}
Sea $A$ un álgebra de Banach. Un \emph{carácter} en $A$ es un funcional lineal multiplicativo no nulo
\[
  \varphi:A\longrightarrow\mathbb{C}.
\]
\end{definicion}

% Fuente: Parte1.pdf, p. 4
\begin{definicion}{1.19}
Sea $A$ un álgebra de Banach conmutativa con unidad. Definimos el \emph{espectro} (o \emph{espacio de caracteres}, o \emph{espectro de Gelfand}) como
\[
  \Omega(A)=\{\text{caracteres }A\to\mathbb{C}\}\subseteq A^*.
\]
Equipamos $\Omega(A)$ con la topología $w^*$.
\end{definicion}

% Fuente: Parte1.pdf, p. 6
\begin{definicion}{1.25}
Sea $A$ un álgebra de Banach con unidad y $x\in A$. Definimos el \emph{radio espectral} como
\[
  r(x)=\max\{\lvert\lambda\rvert:\lambda\in\sigma(x)\}.
\]
\end{definicion}

% Fuente: Parte1.pdf, p. 9
\begin{definicion}{1.27}
Sea $A$ un álgebra de Banach conmutativa con unidad. Se define la \emph{transformada de Gelfand} como el mapeo
\[
  \Lambda:A\longrightarrow C(\Omega(A)),
  \qquad
  x\longmapsto \widehat{x},
  \qquad
  \widehat{x}(\varphi):=\varphi(x).
\]
\end{definicion}

% Fuente: Parte1.pdf, p. 9
\begin{definicion}{1.29}
Un ideal $I$ de $A$ es \emph{modular} si existe $u\in A$ tal que
\[
  a-ua\in I,
  \qquad
  a-au\in I,
  \qquad
  \forall a\in A.
\]
\end{definicion}


\section*{Proposiciones, corolarios, teoremas y otros enunciados}


% Fuente: Parte1.pdf, p. 1
\begin{observacion}{1.2}
Identificamos
\[
  \mathbb{C}\cong\mathbb{C}1_A,
  \qquad
  \lambda\longmapsto\lambda 1_A.
\]
\end{observacion}

% Fuente: Parte1.pdf, p. 1
\begin{observacion}{1.3}
Si $A$ no tiene unidad, es siempre posible agregar una unidad por medio del siguiente procedimiento:
\[
  A^+:=A\times\mathbb{C}
\]
con suma puntual,
\[
  (x,\lambda)\cdot(y,\mu)
  =(xy+\lambda y+\mu x,\lambda\mu),
  \qquad
  \lVert(x,\lambda)\rVert=\lVert x\rVert+|\lambda|.
\]
Entonces $A^+$ es un álgebra de Banach con unidad
\[
  1_{A^+}=(0,1),
\]
y
\[
  x\longmapsto(x,0)
\]
define un encaje isométrico $A\subseteq A^+$.
\end{observacion}

% Fuente: Parte1.pdf, p. 2
\begin{ejemplo}{1.5}
Sea $X$ un espacio localmente compacto y Hausdorff. Definimos el álgebra de funciones acotadas sobre $X$ como
\[
  C_b(X):=\{f:X\to\mathbb{C}: f\text{ continua y acotada}\},
  \qquad
  \lVert f\rVert_\infty=\sup_{x\in X}|f(x)|.
\]
De modo que $C_b(X)$ es un álgebra de Banach. Si además definimos
\[
  f^*(x)=\overline{f(x)},
\]
entonces $C_b(X)$ es un álgebra de Banach involutiva. Si además $X$ es compacto, definimos
\[
  C(X):=C_b(X).
\]
\end{ejemplo}

% Fuente: Parte1.pdf, p. 2
\begin{ejemplo}{1.6}
Si $E$ es un espacio de Banach,
\[
  L(E):=\{T:E\to E:T\text{ lineal y acotado}\}.
\]
$L(E)$ es un álgebra de Banach con la norma de operadores y la composición como producto. $L(E)$ es en general no conmutativa.
\end{ejemplo}

% Fuente: Parte1.pdf, p. 2
\begin{proposicion}{1.8 (Serie de Neumann)}
Sea $A$ un álgebra de Banach con unidad. Si $x\in A$ cumple
\[
  \lVert 1_A-x\rVert<1,
\]
entonces $x$ es invertible y
\[
  x^{-1}=\sum_{n=0}^{\infty}(1_A-x)^n.
\]
\end{proposicion}

% Fuente: Parte1.pdf, p. 2
\begin{observacion}{1.9}
Sea $A$ un álgebra de Banach con unidad y sean $x,z\in A$. Si $z$ es invertible y
\[
  \lVert z-x\rVert<\lVert z^{-1}\rVert^{-1},
\]
entonces $x$ es invertible. Como consecuencia,
\[
  \GL(A)
\]
es un abierto de $A$.
\end{observacion}

% Fuente: Parte1.pdf, p. 3
\begin{corolario}{1.10}
Sea $A$ un álgebra de Banach con unidad. Sean $x\in A$ y $\lambda\in\mathbb{C}$ tales que
\[
  \lVert x\rVert<|\lambda|.
\]
Entonces
\[
  \lambda 1_A-x\in A
\]
es invertible.
\end{corolario}

% Fuente: Parte1.pdf, p. 3
\begin{proposicion}{1.12}
El espectro $\sigma(x)\subseteq\mathbb{C}$ es un subconjunto compacto de
\[
  \{\lambda\in\mathbb{C}:|\lambda|\leq\lVert x\rVert\}
  =B_{\lVert x\rVert}(0).
\]
\end{proposicion}

% Fuente: Parte1.pdf, p. 3
\begin{lema}{1.14}
El resolvente es una función holomorfa, con derivada
\[
  R_x'(z)=-(z1_A-x)^{-2}.
\]
\end{lema}

% Fuente: Parte1.pdf, pp. 3--4
\begin{teorema}{1.15}
El espectro $\sigma(x)$ es no vacío.
\end{teorema}

% Fuente: Parte1.pdf, p. 4
\begin{teorema}{1.16 (Gelfand--Mazur)}
Sea $A$ un álgebra de Banach con unidad en la que todo elemento distinto de cero es invertible, es decir,
\[
  \GL(A)=A\setminus\{0\}.
\]
Entonces
\[
  A=\mathbb{C}1_A.
\]
\end{teorema}

% Fuente: Parte1.pdf, p. 4
\begin{proposicion}{1.18}
Sea $A$ un álgebra de Banach con unidad. Todo carácter
\[
  \varphi:A\to\mathbb{C}
\]
es continuo. De hecho,
\[
  \lVert\varphi\rVert=\varphi(1_A)=1.
\]
\end{proposicion}

% Fuente: Parte1.pdf, pp. 4--5
\begin{proposicion}{1.20}
Sea $A$ un álgebra de Banach con unidad. Supongamos que $J\subseteq A$ es un ideal maximal. Entonces:
\begin{enumerate}[label=(\alph*),leftmargin=2.2em]
  \item $J$ es cerrado;
  \item si $A$ es conmutativa, entonces
  \[
    A/J\cong\mathbb{C}.
  \]
\end{enumerate}
\end{proposicion}

% Fuente: Parte1.pdf, p. 5
\begin{proposicion}{1.21}
Sea $A$ un álgebra de Banach conmutativa con unidad. La asignación
\[
  \Omega(A)\longrightarrow\{\text{ideales maximales}\},
  \qquad
  \varphi\longmapsto\ker\varphi,
\]
es una biyección.
\end{proposicion}

% Fuente: Parte1.pdf, p. 5
\begin{teorema}{1.22}
Sea $A$ un álgebra de Banach conmutativa con unidad y $x\in A$. Entonces
\[
  \sigma(x)=\{\varphi(x):\varphi\in\Omega(A)\}.
\]
\end{teorema}

% Fuente: Parte1.pdf, pp. 5--6
\begin{proposicion}{1.23}
Sea $A$ un álgebra de Banach con unidad y $x,y\in A$. Entonces
\[
  1_A-xy\in\GL(A)
  \quad\Longleftrightarrow\quad
  1_A-yx\in\GL(A).
\]
\end{proposicion}

% Fuente: Parte1.pdf, p. 6
\begin{corolario}{1.24}
Sea $A$ un álgebra de Banach con unidad. Si $x,y\in A$, entonces
\[
  \sigma(xy)\cup\{0\}
  =
  \sigma(yx)\cup\{0\}.
\]
\end{corolario}

% Fuente: Parte1.pdf, pp. 6--8
\begin{teorema}{1.26 (Fórmula del Radio Espectral)}
Sea $A$ un álgebra de Banach con unidad y $x\in A$. Entonces el límite
\[
  \lim_{n\to\infty}\lVert x^n\rVert^{1/n}
\]
existe y
\[
  r(x)
  =\lim_{n\to\infty}\lVert x^n\rVert^{1/n}
  =\inf_{n\in\mathbb{N}}\lVert x^n\rVert^{1/n}.
\]
\end{teorema}

% Fuente: Parte1.pdf, p. 9
\begin{observacion}{1.28}
La transformada de Gelfand es lineal y multiplicativa:
\[
  \widehat{x+\lambda y}(\varphi)
  =\varphi(x+\lambda y)
  =\varphi(x)+\lambda\varphi(y)
  =\widehat{x}(\varphi)+\lambda\widehat{y}(\varphi),
\]
\[
  \widehat{xy}(\varphi)
  =\varphi(xy)
  =\varphi(x)\varphi(y)
  =\widehat{x}(\varphi)\widehat{y}(\varphi).
\]
Además,
\[
  \lVert\widehat{x}\rVert
  =\sup_{\varphi\in\Omega(A)}|\varphi(x)|
  =r(x)
  \leq\lVert x\rVert.
\]
Por lo tanto, la transformada de Gelfand es un homomorfismo contractivo que preserva la unidad.
\end{observacion}

% Fuente: Parte1.pdf, p. 9
\begin{ejemplo}{1.30}
Sea $X$ localmente compacto y Hausdorff, y sea $x\in X$. Definimos
\[
  I_x=\{f\in C_0(X):f(x)=0\}.
\]
Entonces $I_x$ es un ideal modular. Además, $I_x$ tiene codimensión $1$ en $C_0(X)$ y, por tanto, es maximal.
\end{ejemplo}

% Fuente: Parte1.pdf, p. 9
\begin{proposicion}{1.31}
Un ideal $I$ es modular si y sólo si $A/I$ tiene unidad.
\end{proposicion}


\end{document}
